% Generated draft: requires mathematical and source review.
\documentclass[11pt]{article}
\usepackage[margin=25mm]{geometry}
\usepackage{fontspec}
\setmainfont{Latin Modern Roman}
\usepackage{amsmath,amssymb,mathtools,unicode-math}
\setmathfont{Latin Modern Math}
\usepackage{enumitem,needspace,xurl,hyperref}
\setlength{\parindent}{0pt}
\setlength{\parskip}{5pt}
\setlength{\emergencystretch}{3em}
\newcommand{\N}{\mathbb{N}}
\newcommand{\Z}{\mathbb{Z}}
\newcommand{\Q}{\mathbb{Q}}
\newcommand{\R}{\mathbb{R}}
\newcommand{\C}{\mathbb{C}}
\newcommand{\F}{\mathbb{F}}
\newcommand{\A}{\mathbb{A}}
\newcommand{\PP}{\mathbb{P}}
\newcommand{\E}{\mathbb{E}}
\newcommand{\eps}{\varepsilon}
\newcommand{\dd}{\,\mathrm{d}}
\begin{document}
\section*{UnsolvedMath ID: 3100067}
\textbf{Source problem code:} catalogue entry
\subsection{Definitions and mathematical statement}
Let $p_n$ be the $n$th prime, in increasing order, and let $g_n=p_{n+1}-p_n$ be the gap between consecutive primes. The supplied question asks how fast these gaps grow, and whether they grow more slowly than $n^{\varepsilon}$ for every $\varepsilon>0$. One precise eventual-inequality reading is
\[
 \forall\varepsilon>0\ \exists N_{\varepsilon}\ \forall n\ge N_{\varepsilon},
 \qquad g_n<n^{\varepsilon}.
\]
The original exponent appears as the corrupted character ``ľ'', and ``slower than'' does not explicitly specify an eventual inequality, a big-$O$ bound, or a little-$o$ bound. The displayed assertion states an interpretation for review rather than silently resolving the wording.

The research summary instead measures the gap by the prime itself and discusses
\[
 \forall\varepsilon>0\ \exists N_{\varepsilon}\ \forall n\ge N_{\varepsilon},
 \qquad g_n<p_n^{\varepsilon},
\]
and the closely related interval formulation that, for each $\varepsilon>0$, every sufficiently large real $x$ has a prime in $(x,x+x^{\varepsilon})$. The original use of $n$ is not silently replaced by $p_n$.

As supplied literature context, the record attributes a prime-in-short-interval bound with exponent $0.525$ to Baker--Harman--Pintz (2001), describes the unconditional gap exponent as approximately $0.525$, and reports under the Riemann hypothesis
\[
 g_n=O\bigl(\sqrt{p_n}(\log p_n)^2\bigr).
\]
The same summary says that the assertion for every positive $\varepsilon$ ``follows from RH.'' That inference conflicts with the displayed RH estimate, which only supplies a square-root-scale bound up to logarithms and does not establish arbitrarily small positive exponents. The all-$\varepsilon$ implication is therefore retained as an unresolved source claim, not asserted as a theorem.
\subsection{Short English statement}
Do consecutive-prime gaps eventually lie below every positive power of their index, or, in the research summary's formulation, below every positive power of the prime? The record labels the question open with partial progress. Its cited unconditional exponent is $0.525$, and its stated RH bound has exponent $1/2$ up to logarithms. Its separate claim that RH implies the full all-positive-exponents conjecture requires correction or justification.

The source page's open label belongs to its October 2020 snapshot; provenance explicitly says later resolution was not checked item by item. The later registry summary is supplied context, not a newly verified literature survey.
\subsection{Sources}
\begin{itemize}
\item Local UnsolvedMath registry, permanent ID 3100067, source code catalogue entry, ``How quickly do the gaps between successive primes grow.'' Dataset curation and metadata: Ulam AI and UnsolvedMath Contributors, CC BY 4.0.
\item Cooper, \emph{Combinatorial Problems I Like}, October 2020 snapshot, Classical Number Theory, source-order bullet 67, \url{https://people.math.sc.edu/cooper/combprob.html}; HTML extraction, accessed 2026-07-29. The source's ``See this'' link has no target in the supplied record.
\item Research-summary attribution: Baker--Harman--Pintz (2001), short prime intervals with exponent $0.525$. No complete bibliographic entry or locator for that result or the stated RH estimate is supplied.
\item Provenance marks status evidence NEEDS\_REVIEW and rights NEEDS\_REVIEW: public access was verified, but redistribution permission was not established. The source-specific rights qualification is retained separately from the dataset licence.
\end{itemize}
\end{document}
